\par\SourceLeaf{305}
\indent Überhaupt verhält es sich mit allen solchen neuen Calculs so, dass man durch sie
nichts leisten kann, was nicht auch ohne sie zu leisten wäre; der Vortheil ist aber
der, dass, wenn ein solcher Calcul dem innersten Wesen vielfach vorkommender
Bedürfnisse correspondirt, jeder, der sich ihn ganz angeeignet hat, auch ohne die
gleichsam unbewussten Inspirationen des Genies, die niemand \textls{erzwingen} kann, die dahin
gehörigen Aufgaben lösen, ja selbst in so verwickelten Fällen gleichsam mechanisch
lösen kann, wo ohne eine solche Hülfe auch das Genie ohnmächtig wird. So ist es mit
der Erfindung der Buchstabenrechnung überhaupt; so mit der Differentialrechnung
gewesen; so ist es auch (wenn auch in partielleren Sphären) mit \textsc{Lagranges}
Variationsrechnung, mit meiner Congruenzenrechnung und mit \textsc{Möbius}' Calcul. Es
werden durch solche Conceptionen unzählige Aufgaben, die sonst vereinzelt stehen, und
jedesmal neue Efforts (kleinere oder grössere) des Erfindungsgeistes erfordern,
gleichsam zu einem organischen Reiche. \mbox{.\,.\,.\,.\,.\,.\,.\,.\,.}

\Needspace{9\baselineskip}\begin{center}
\medskip

[II.]

\textbf{[DER BARYCENTRISCHE CALCUL UND DER RESULTANTENCALCUL.]}

[1.]
\end{center}

Der barycentrische Calcul findet sein Gegenstück in einem andern (vermuthlich noch
umfassendern) Calcul, den man den Resultantencalcul nennen könnte. So wie der erste
sich mit Punkten beschäftigt, in denen man schwere Massen voraussetzt; so würde
letzterer zum Gegenstande haben Linien, in welchen Kräfte wirken. Sind $a$, $b$, $c$,
$d$ u. s. w. solche Linien, in denen, in jeder in bestimmtem Sinn, Kräfte wirken, die
den Zahlen $\alpha$, $\beta$, $\gamma$, $\delta$ u. s. w. proportional sind, so würde
die Gleichung

\[
\alpha a + \beta b + \gamma c + \delta d + \text{u. s. w.} = 0
\]
\noindent bedeuten, dass diese Kräfte einander das Gleichgewicht halten.
